\section{Running-time analysis and stored timing evidence}
Let $n$ be clinics, $\mu$ GA population, $k$ candidate-list size and $L$ maximum route length.
These describe the cost of a step, not how many steps a good answer needs.

\textbf{Clarke--Wright.} There are $O(n^2)$ savings. Sorting costs $O(n^2\log n)$. At most
$n-1$ successful merges each touch $O(n)$ clinics in this implementation. This adds $O(n^2)$
and does not dominate sorting.

\textbf{GA generation.} A generation attempts $\mu$ offspring. Per child, crossover, mutation,
fast unlimited Split and hashing cost at most $O(n)$ each. Reordering after replacement costs
$O(\mu)$. The ordinary generation cost is
\[
 O\bigl(\mu(n+\mu)\bigr).
\]
If every child needs fleet-limited fast Split, a conservative bound is $O(\mu(Kn+\mu))$.
With ordinary Bellman Split, replace the decoding term by $O(nL)$ per child. Keeping $\mu$
explicit matters: population 100 is larger than some main benchmarks.

\textbf{TS iteration.} The move generator checks $O(nk)$ candidates. Each uses constant-time
delta evaluation and a constant-size tabu check. Applying the selected move costs at most
$O(n)$, so an iteration costs $O(nk)$: linear for fixed $k$, quadratic for $k=n-1$.

\textbf{Delta evaluation.} Fully recomputing a solution costs $O(n)$ per candidate, giving
$O(n^2k)$ per iteration instead of $O(nk)$. Both solvers store an $O(n^2)$ distance matrix.
GA adds $O(\mu n)$ population memory; TS adds an $O(n^2)$ tabu matrix.

\begin{figure}[h]\centering
\includegraphics[width=\linewidth]{\figdir/complexity.pdf}
\caption{Saved timing data on log--log axes. Left: time per step, with fitted slopes in
parentheses. Right: time per candidate for delta and full evaluation.}
\end{figure}

\textbf{Measurements.} The stored study uses synthetic sizes 25--3200, three repetitions, fixed
iteration counts and approximately constant route lengths. Timing two run lengths and subtracting
removes initialization. Its GA population is 50, rather than the main comparison's 100.
Slopes are 1.08 for fixed-list TS (fit through $n=1600$), 2.12 for full-list TS and 0.84 for GA.
These broadly match predictions 1, 2 and 1 for fixed population and route length. The plotted
$n=3200$ TS point is outside its fit; the sharp increase is consistent with cache and memory
effects. GA's fixed overhead matters more for small sizes. A measured slope supports the analysis
over those sizes, but is not a proof of asymptotic complexity.
